\begin{proof} 1. \\
  \begin{itemize}
      \item[]\hspace{-0.5cm} ($(1) \rightarrow (2)$)\\
        Prop7.1.5.1 より明らか.
      \item[]\hspace{-0.5cm} ($(2) \rightarrow (1)$)\\
        対偶を示す.\\
        仮定とProp7.1.5.2より,
        $\BijMap {\seigen {C_A}{A}}{A}{\nat}$.\\
        Def7.1.4.2より$A$は可算無限集合である.\\
        よって対偶は真.
  \end{itemize}
\end{proof}


\begin{proof} 2. \\
  まず補題として以下の証明を行う.
  \begin{equation*}
   A \neq \emptyset \rightarrow \forall m,n \in A(C_A(m) \leq C_A(n) \leftrightarrow m \leq n )
  \end{equation*}
  $A \neq \emptyset$を仮定する.\\
  \begin{itemize}
      \item[]\hspace{-0.5cm} ($(1) \rightarrow (2)$)\\
        $C_A$の定義より明らか.
      \item[]\hspace{-0.5cm} ($(2) \rightarrow (1)$)\\
        任意に$m,n \in A$をとり,
        $C_A(m) \leq C_A(n)$とする.\\
        また,
        $(1) \rightarrow (2)$と,
        $\seigen {C_A}{A}$は単射より以下は真.
          \begin{equation}
            \forall a,b \in A(a < b \rightarrow C_A(a) < C_A(b))
          \end{equation}
        命題2.9.1より,
        $\nat$は大小関係に関して全順序集合だから,
        以下は真.
        \begin{equation*}
          \forall a, b \in A (a \leq b \lor b \leq a )
        \end{equation*}
        これと(7.1)の対偶を合わせると,
        以下は真.
        \begin{equation}
            \forall a,b \in A(C_A(b) \geq C_A(a) \rightarrow b \geq a)
        \end{equation}
        $m,n$を(7.2)に代入すると$C_A(m) \leq C_A(n)$より,
        $m \leq n$は真.
  \end{itemize}
  次に2を示す.\\
  goal :$\Leftrightarrow A \neq \emptyset \rightarrow
  \exists l (\seigen{C_A}{A}(l) = 0 \land \forall n \in A(l \leq n))$\\
  $A \neq \emptyset$を仮定する.\\
  $l := \seigen{C_A}{A}^{-1}(0)$とおくと$\seigen{C_A}{A}(l) = 0$.\\
  いま,
  $\forall n \in \nat (0 \leq x)$より,
  $\forall n \in A (C_A(l) \leq C_A(n))$は真.\\
  これと補題を合わせると,
  $\forall n \in A (l \leq n)$.\\
  そんな$l$の存在よりgoalは真.
\end{proof}
